% ============================================================
% Cheat sheet, argomento: il pattern Builder (dispensa, capitolo 15)
% ============================================================
\argomento{Builder}{Dispensa: capitolo 15. Pattern creazionale della GoF}[Pattern Builder]

\scheda{Carta d'identità}{
  \pintento{Rendere la costruzione di un oggetto complesso \textbf{flessibile e scomponibile in passi}, separandola
    dal costruttore dell'oggetto; con un Director, gli \textbf{stessi passi} producono \textbf{rappresentazioni
    diverse}.}
  \ptesto{«molti parametri, alcuni obbligatori e molti facoltativi», «configurazioni diverse dello stesso
    prodotto», «costruito un pezzo alla volta», «in futuro le opzioni aumenteranno», «controllare i dati prima di
    creare l'oggetto», «la stessa procedura deve produrre risultati diversi» (una casa e un grattacielo, una scheda
    tecnica e un preventivo).}
  \pcodice{una classe \code{XBuilder} con gli stessi campi del prodotto; i parametri obbligatori nel suo costruttore;
    metodi \code{withY(\ldots)} (o \code{setY}) che impostano un'opzione e fanno \textbf{\code{return this}}, così le
    chiamate si concatenano; un metodo finale \textbf{\code{build()}} (o \code{getResult()}) che controlla i dati e
    fa \code{new Prodotto(\ldots)}. Il prodotto ha campi \code{final} e un costruttore non pubblico.}
  \pdiagramma{il Builder con una dipendenza «create» verso il prodotto, e operazioni che restituiscono il Builder
    stesso. Con il Director: un'interfaccia (o classe astratta) Builder con più ConcreteBuilder, e il Director che la
    usa.}
  \pvarianti{senza Director (il client chiama i passi: \vv{builder-esempio}) o con il Director che conosce le
    «ricette» (\vv{builder-director}); il builder come classe annidata statica del prodotto
    (\code{Computer.builder()}) \oltre. Nella libreria di Java: \code{StringBuilder}, \code{Locale.Builder},
    \code{DateTimeFormatterBuilder}, \code{HttpRequest.newBuilder()}.}
  \pconfondere{un costruttore con dieci parametri senza nome o tanti \code{set} su un oggetto già creato (che resta
    a metà e modificabile per sempre); il Decorator, dove le catene sono di oggetti uno dentro l'altro
    (\code{new A(new B(\ldots))}), non di chiamate sullo stesso oggetto.}
}

\begin{struttura}{builder-struttura}{La struttura}
\begin{tikzpicture}
  \umlclass[2.2cm]{d}{(0,0)}{Director}{}{+ construct()}
  \umlclass[2.4cm]{b}{(3.9,0)}{\textit{Builder}}{}{+ buildPart()}
  \umlclass[2.6cm]{c}{(3.9,-2.4)}{ConcreteBuilder}{}{+ buildPart()\\+ getResult()}
  \node[classe1=1.6cm] (p) at (7.3,-2.4) {Product};
  \draw[aggr] (d.east) -- (b.west);
  \estO{(b.west)}{}{builder}
  \draw[gen] (c.north) -- (b.south);
  \draw[dep] (c.east) -- node[stereo, above]{«create»} (p.west);
  \node[notauml, anchor=north, text width=2.6cm] (n) at (0,-1.5) {for all parti:\\\hspace*{1em}builder.buildPart()};
  \draw[tratt] (d.south) -- (n.north);
\end{tikzpicture}
\tcblower
\textbf{Builder}: l'interfaccia dei passi di costruzione. \textbf{ConcreteBuilder}: li implementa, tiene il
prodotto in costruzione e lo consegna alla fine (\code{getResult}, \code{build}). \textbf{Director}: conosce la
sequenza dei passi e la ripete con qualunque builder. \textbf{Product}: l'oggetto complesso. \code{getResult()} sta
solo nei builder concreti, perché ogni builder può produrre un tipo diverso.
\end{struttura}

\begin{programma}{builder-esempio}{L'esempio del corso: costruire un computer (senza Director)}
\begin{lstlisting}[style=javacs, name=Main.java]
public class Main {
    public static void main(String[] args) {
        Computer gaming =
            new ComputerBuilder("AMD Ryzen 7", 32)
                .withDedicatedGPU("RX 7800")
                .withSSD(2000)
                .includeOS()
                .build();
        System.out.println(gaming);
        try {
            new ComputerBuilder("Intel Core i3", 2)
                .build();
        } catch (IllegalStateException e) {
            System.out.println(e.getMessage());
        }
    }
}
class Computer {
    private final String cpu, gpu, disco;
    private final int ramGB;
    private final boolean hasOS;
    Computer(String cpu, int ramGB, String gpu,
             String disco, boolean hasOS) {
        this.cpu = cpu;     this.ramGB = ramGB;
        this.gpu = gpu;     this.disco = disco;
        this.hasOS = hasOS;
    }
    public String toString() {
        String os = hasOS ? ", con OS" : "";
        return cpu + ", " + ramGB + " GB, " + gpu
                + ", " + disco + os;
    }
}
class ComputerBuilder {
    private final String cpu;      // obbligatori
    private final int ramGB;
    private String gpu = "GPU integrata"; // opzioni
    private String disco = "nessun disco";
    private boolean hasOS = false;

    ComputerBuilder(String cpu, int ramGB) {
        this.cpu = cpu;
        this.ramGB = ramGB;
    }
    ComputerBuilder withDedicatedGPU(String gpu) {
        this.gpu = gpu;
        return this;
    }
    ComputerBuilder withSSD(int gb) {
        disco = "SSD " + gb + " GB";
        return this;
    }
    ComputerBuilder includeOS() {
        hasOS = true;
        return this;
    }
    Computer build() {
        if (ramGB < 4)
            throw new IllegalStateException(
                    "RAM insufficiente: min 4 GB");
        return new Computer(cpu, ramGB, gpu, disco,
                            hasOS);
    }
}
\end{lstlisting}
\tcblower
\begin{tikzpicture}
  \umlclass[3.2cm]{m}{(0,2.4)}{Main}{}{\underline{+ main(args: String [*])}}
  \umlclass[6.4cm]{b}{(0,-0.5)}{ComputerBuilder}{- cpu: String \{readOnly\}\\- ramGB: int \{readOnly\}\\- gpu: String = "GPU integrata"\\- disco: String = "nessun disco"\\- hasOS: boolean = false}{\textasciitilde\ ComputerBuilder(cpu: String, ramGB: int)\\\textasciitilde\ withDedicatedGPU(gpu: String): ComputerBuilder\\\textasciitilde\ withSSD(gb: int): ComputerBuilder\\\textasciitilde\ includeOS(): ComputerBuilder\\\textasciitilde\ build(): Computer}
  \umlclass[3.6cm]{c}{(0,-5.1)}{Computer}{- cpu: String \{readOnly\}\\- gpu: String \{readOnly\}\\- disco: String \{readOnly\}\\- ramGB: int \{readOnly\}\\- hasOS: boolean \{readOnly\}}{\textasciitilde\ Computer(\ldots)\\+ toString(): String}
  \draw[dep] (m.south) -- node[stereo, right]{«use»} (b.north);
  \draw[dep] (b.south) -- node[stereo, right]{«create»} (c.north);
\end{tikzpicture}

\medskip
\begin{lstlisting}[style=shellcs, name=Terminale]
$ javac Main.java
$ java Main
AMD Ryzen 7, 32 GB, RX 7800, SSD 2000 GB, con OS
RAM insufficiente: min 4 GB
\end{lstlisting}
\regola{Ogni \code{with\ldots} restituisce \code{this}: la catena lavora sempre sullo \textbf{stesso} builder.
\code{build()} controlla i dati e solo alla fine crea il \code{Computer}, che non cambia più (campi \code{final}).
Il secondo computer non viene creato: la regola sta nel builder, non nel prodotto.}
\end{programma}

\begin{programma}{builder-director}{Con il Director: stessa ricetta, prodotti diversi}
\begin{lstlisting}[style=javacs, name=Negozio.java]
public class Negozio {
    public static void main(String[] args) {
        Configuratore dir = new Configuratore();
        SchedaTecnica scheda = new SchedaTecnica();
        dir.pcGaming(scheda);      // stessa ricetta
        System.out.print(scheda.getResult());
        Preventivo prev = new Preventivo();
        dir.pcGaming(prev);        // altro prodotto
        System.out.println("Prezzo: "
                + prev.getResult());
    }
}
interface BuilderPC {                    // Builder
    void cpu(String modello, int ramGB);
    void gpu(String modello);
    void ssd(int capacitaGB);
}
class Configuratore {                    // Director
    void pcGaming(BuilderPC b) {
        b.cpu("AMD Ryzen 7", 32);
        b.gpu("AMD Radeon RX 7800");
        b.ssd(2000);
    }
    void pcUfficio(BuilderPC b) {
        b.cpu("Intel Core i3", 8);
        b.ssd(256);
    }
}
class SchedaTecnica implements BuilderPC {
    private final StringBuilder testo =
            new StringBuilder();
    public void cpu(String modello, int ramGB) {
        testo.append("CPU " + modello + ", RAM "
                + ramGB + " GB\n");
    }
    public void gpu(String modello) {
        testo.append("GPU " + modello + "\n");
    }
    public void ssd(int gb) {
        testo.append("SSD " + gb + " GB\n");
    }
    String getResult() { return testo.toString(); }
}
class Preventivo implements BuilderPC {
    private double totale = 0;
    public void cpu(String modello, int ramGB) {
        totale += 250 + 4 * ramGB;
    }
    public void gpu(String m) { totale += 550; }
    public void ssd(int gb) { totale += 0.08 * gb; }
    double getResult() { return totale; }
}
\end{lstlisting}
\tcblower
\begin{tikzpicture}
  \umlclass[3.4cm]{d}{(0,0)}{Configuratore}{}{\textasciitilde\ pcGaming(b: BuilderPC)\\\textasciitilde\ pcUfficio(b: BuilderPC)}
  \umlclass[3.6cm]{i}{(4.1,0)}{\stereotipo{interface}{BuilderPC}}{}{+ cpu(modello: String, ramGB: int)\\+ gpu(modello: String)\\+ ssd(capacitaGB: int)}
  \umlclass[3.0cm]{s}{(2.2,-2.9)}{SchedaTecnica}{- testo: StringBuilder}{\textasciitilde\ getResult(): String}
  \umlclass[2.7cm]{p}{(5.4,-2.9)}{Preventivo}{- totale: double = 0}{\textasciitilde\ getResult(): double}
  \draw[dep] (d.east) -- node[stereo, above]{«use»} (d.east -| i.west);
  \draw[realizza] (s.north) |- ($(i.south)+(0,-0.45)$) -- (i.south);
  \draw[realizza] (p.north) |- ($(i.south)+(0,-0.45)$);
\end{tikzpicture}

\medskip
\begin{lstlisting}[style=shellcs, name=Terminale]
$ javac Negozio.java
$ java Negozio
CPU AMD Ryzen 7, RAM 32 GB
GPU AMD Radeon RX 7800
SSD 2000 GB
Prezzo: 1088.0
\end{lstlisting}
\regola{Il \textbf{Director} (\code{Configuratore}) conosce le ricette ma non il prodotto; i due \textbf{ConcreteBuilder}
eseguono gli stessi passi e producono una scheda (\code{String}) e un prezzo (\code{double}): per questo
\code{getResult()} non sta nell'interfaccia. Prezzi di fantasia: $378 + 550 + 160 = 1088$.}
\end{programma}
